\documentclass[10pt, a4paper]{report}
\usepackage[utf8]{inputenc}
\usepackage{geometry}
\usepackage{longtable}
\usepackage{array} % For better column definitions in longtable
\usepackage{enumitem} % Required for itemize and enumerate options
\usepackage{helvet} % Using a sans-serif font for better readability
\renewcommand{\familydefault}{\sfdefault}

\geometry{left=0.5cm, right=0.5cm, top=2cm, bottom=2cm}

\title{Non-Functional Test Cases: Performance Testing \\ \large Library Management System}
\author{Test Team (Coders \& Non-Coders/Specification Experts)}
\date{June 5, 2025}

\begin{document}
\maketitle
\begin{center}
    \textbf{Note:} These test cases are designed to evaluate the system's performance characteristics, including responsiveness, speed, and stability under various conditions. Metrics are illustrative and should be benchmarked against specific project requirements. Refer to the SRS and Test Plan for complete context.
\end{center}

\section*{IV.C Non-Functional Test Cases: Performance Testing}

\begin{longtable}{|>{\raggedright\arraybackslash}p{3.5cm}|>{\raggedright\arraybackslash}p{4cm}|>{\raggedright\arraybackslash}p{5cm}|>{\raggedright\arraybackslash}p{5.5cm}|}
\caption{Performance Test Cases}\label{tab:perf_test_cases}\\
\hline
\textbf{Test Case ID} & \textbf{Test Title} & \textbf{Test Steps} & \textbf{Expected Results} \\
\hline
\endfirsthead
\multicolumn{4}{c}%
{{\bfseries \tablename\ \thetable{} -- continued from previous page}} \\
\hline
\textbf{Test Case ID} & \textbf{Test Title} & \textbf{Test Steps} & \textbf{Expected Results} \\
\hline
\endhead
\hline \multicolumn{4}{|r|}{{Continued on next page}} \\ \hline
\endfoot
\hline
\endlastfoot

% UI Load & Responsiveness Testing (Non-Coder Focus)
\multicolumn{4}{|l|}{\textbf{Sub-Section: UI Load \& Responsiveness Testing}} \\ \hline
TC\_PERF\_UI\_001 & Page load time for book list with few records & 
    \begin{itemize}[noitemsep, topsep=0pt, partopsep=0pt, leftmargin=*]
        \item \textbf{Precondition:} Database contains a small number of records (e.g., < 50 books).
        \item Login as a Librarian or Member.
        \item Navigate to the book list page ('/books').
        \item Measure the time from navigation start to when the page is fully rendered and interactive.
    \end{itemize} & 
    \begin{itemize}[noitemsep, topsep=0pt, partopsep=0pt, leftmargin=*]
        \item The page load time should be within the acceptable baseline (e.g., < 2 seconds).
        \item The UI is responsive and usable immediately after rendering.
    \end{itemize} \\ \hline
TC\_PERF\_UI\_002 & Page load time for book list with many records (Volume) & 
    \begin{itemize}[noitemsep, topsep=0pt, partopsep=0pt, leftmargin=*]
        \item \textbf{Precondition:} Database contains a large number of records (e.g., > 5,000 books).
        \item Login as a Librarian or Member.
        \item Navigate to the book list page ('/books').
        \item Measure the time until the page is fully rendered and interactive.
    \end{itemize} & 
    \begin{itemize}[noitemsep, topsep=0pt, partopsep=0pt, leftmargin=*]
        \item The page load time meets the defined threshold for large data sets (e.g., < 5 seconds).
        \item UI remains responsive; features like sorting or filtering might show a loading indicator but do not freeze the application.
    \end{itemize} \\ \hline
TC\_PERF\_UI\_003 & UI responsiveness during data entry on a high-load page & 
    \begin{itemize}[noitemsep, topsep=0pt, partopsep=0pt, leftmargin=*]
        \item Navigate to a data-heavy page (e.g., Librarian Dashboard with many widgets).
        \item While the page is rendering or widgets are loading data, attempt to interact with a simple UI element (e.g., open a dropdown menu).
    \end{itemize} & 
    \begin{itemize}[noitemsep, topsep=0pt, partopsep=0pt, leftmargin=*]
        \item User interactions are not blocked.
        \item The UI responds to input within a reasonable time (e.g., < 500ms) without noticeable lag.
    \end{itemize} \\ \hline

% API Response Time Testing (Coder Focus)
\multicolumn{4}{|l|}{\textbf{Sub-Section: API Response Time Testing}} \\ \hline
TC\_PERF\_API\_001 & API response time for fetching all books (Normal Load) & 
    \begin{itemize}[noitemsep, topsep=0pt, partopsep=0pt, leftmargin=*]
        \item \textbf{Precondition:} Database contains a moderate number of book records (e.g., 500).
        \item Using an API testing tool (e.g., Postman), send a GET request to the `/api/book/getAll` endpoint.
        \item Measure the server response time.
    \end{itemize} & 
    \begin{itemize}[noitemsep, topsep=0pt, partopsep=0pt, leftmargin=*]
        \item The API response time is within the defined benchmark (e.g., < 300ms).
        \item The HTTP status code is 200 OK.
    \end{itemize} \\ \hline
TC\_PERF\_API\_002 & API response time for creating a new entity (e.g., Borrowal) & 
    \begin{itemize}[noitemsep, topsep=0pt, partopsep=0pt, leftmargin=*]
        \item Using an API testing tool, send a POST request with a valid payload to the `/api/borrowal/add` endpoint.
        \item Measure the server response time.
    \end{itemize} & 
    \begin{itemize}[noitemsep, topsep=0pt, partopsep=0pt, leftmargin=*]
        \item The API response time for the creation operation is within the benchmark (e.g., < 400ms).
        \item The HTTP status code is 201 Created (or similar success code).
    \end{itemize} \\ \hline
TC\_PERF\_API\_003 & API response time for user authentication & 
    \begin{itemize}[noitemsep, topsep=0pt, partopsep=0pt, leftmargin=*]
        \item Using an API testing tool, send a POST request with valid user credentials to the `/api/auth/login` endpoint.
        \item Measure the server response time.
    \end{itemize} & 
    \begin{itemize}[noitemsep, topsep=0pt, partopsep=0pt, leftmargin=*]
        \item Authentication is a critical path; response time should be fast (e.g., < 250ms).
        \item The HTTP status code is 200 OK.
    \end{itemize} \\ \hline

% Volume Testing
\multicolumn{4}{|l|}{\textbf{Sub-Section: Volume Testing}} \\ \hline
TC\_PERF\_VOL\_001 & System performance with large data volume in multiple tables & 
    \begin{itemize}[noitemsep, topsep=0pt, partopsep=0pt, leftmargin=*]
        \item \textbf{Precondition:} Populate the database with an abnormally large volume of data:
        \begin{itemize}
            \item 20,000+ Books
            \item 5,000+ Users
            \item 50,000+ Borrowal records
        \end{itemize}
        \item Perform a series of common user actions (e.g., login, search for a book, view user list, view borrowal history).
    \end{itemize} & 
    \begin{itemize}[noitemsep, topsep=0pt, partopsep=0pt, leftmargin=*]
        \item The system remains stable and does not crash or become unresponsive.
        \item Data retrieval and display operations complete within defined acceptable limits, though slower than with few records.
        \item Database queries do not time out.
    \end{itemize} \\ \hline

% Stress Testing
\multicolumn{4}{|l|}{\textbf{Sub-Section: Stress Testing}} \\ \hline
TC\_PERF\_STRS\_001 & System stability under high concurrent user load & 
    \begin{itemize}[noitemsep, topsep=0pt, partopsep=0pt, leftmargin=*]
        \item \textbf{Precondition:} Use a load testing tool (e.g., JMeter, Gatling).
        \item Simulate a high number of concurrent users (e.g., 100 users) performing actions simultaneously for a sustained period (e.g., 15 minutes). Actions include:
        \begin{itemize}
            \item Logging in
            \item Fetching book lists
            \item Adding new borrowals
        \end{itemize}
    \end{itemize} & 
    \begin{itemize}[noitemsep, topsep=0pt, partopsep=0pt, leftmargin=*]
        \item The system remains operational and does not crash.
        \item The average API response time degradation is within an acceptable percentage (e.g., does not increase by more than 50% from the baseline).
        \item The error rate (e.g., HTTP 5xx errors) is below a minimal threshold (e.g., < 1%).
    \end{itemize} \\ \hline
TC\_PERF\_STRS\_002 & System recovery after a stress period & 
    \begin{itemize}[noitemsep, topsep=0pt, partopsep=0pt, leftmargin=*]
        \item Following the completion of the TC\_PERF\_STRS\_001 test, stop the load generation.
        \item Wait for a brief period (e.g., 2 minutes).
        \item Perform a single-user check by navigating key pages (Login, Books, Dashboard).
    \end{itemize} & 
    \begin{itemize}[noitemsep, topsep=0pt, partopsep=0pt, leftmargin=*]
        \item The system recovers and returns to normal performance levels.
        \item Response times for the single user should return to the baseline measurements recorded in the API and UI tests.
        \item No data corruption has occurred.
    \end{itemize} \\ \hline

\end{longtable}

\end{document}